% Generated draft: requires mathematical and source review.
\documentclass[11pt]{article}
\usepackage[margin=25mm]{geometry}
\usepackage{fontspec}
\setmainfont{Latin Modern Roman}
\usepackage{amsmath,amssymb,mathtools,unicode-math}
\setmathfont{Latin Modern Math}
\usepackage{enumitem,needspace,xurl,hyperref}
\setlength{\parindent}{0pt}
\setlength{\parskip}{5pt}
\setlength{\emergencystretch}{3em}
\newcommand{\N}{\mathbb{N}}
\newcommand{\Z}{\mathbb{Z}}
\newcommand{\Q}{\mathbb{Q}}
\newcommand{\R}{\mathbb{R}}
\newcommand{\C}{\mathbb{C}}
\newcommand{\F}{\mathbb{F}}
\newcommand{\A}{\mathbb{A}}
\newcommand{\PP}{\mathbb{P}}
\newcommand{\E}{\mathbb{E}}
\newcommand{\eps}{\varepsilon}
\newcommand{\dd}{\,\mathrm{d}}
\begin{document}
\section*{UnsolvedMath ID: 701296}
\textbf{Source problem code:} catalogue entry
\subsection{Definitions and mathematical statement}
A group \(G\) is locally graded if every nontrivial finitely generated subgroup \(H\leq G\) has a proper subgroup \(K<H\) with \([H:K]<\infty\). A group \(P\) is polycyclic if it has a finite subnormal series
\[
1=P_0\trianglelefteq P_1\trianglelefteq\cdots\trianglelefteq P_n=P
\]
whose factors \(P_i/P_{i-1}\) are cyclic. Here, almost polycyclic means having a polycyclic subgroup of finite index, and almost soluble means having a soluble subgroup of finite index. A group is soluble if its derived series reaches the trivial subgroup after finitely many steps. The maximal condition and minimal condition mean, respectively, that every ascending chain and every descending chain of subgroups stabilizes.

The question is whether, for every locally graded group \(G\) and all subgroups \(A,B\leq G\),
\[
\left(G=AB\ \text{and}\ A,B\ \text{are almost polycyclic}\right)
\quad\Longrightarrow\quad
G\ \text{is almost polycyclic}.
\]
Here \(AB=\{ab:a\in A,\ b\in B\}\), so the hypothesis requires a setwise product decomposition. The source equivalently describes each of the two factors as almost soluble and satisfying the maximal condition on subgroups.

For the supplied comparison, a Chernikov group means a group having a normal subgroup of finite index that is a direct product of finitely many quasicyclic groups \(C_{p^\infty}\), where \(C_{p^\infty}\) is the union of an ascending sequence of cyclic groups of orders \(p,p^2,p^3,\ldots\); the empty direct product is allowed. The source attributes to N. S. Chernikov the assertion that every locally graded group which is a setwise product of two Chernikov subgroups is Chernikov. It equivalently describes those factors as almost soluble subgroups satisfying the minimal condition. This comparison is supplied background, not an independently verified result in this draft.
\subsection{Short English statement}
Must a locally graded group that is the setwise product of two subgroups, each having a polycyclic subgroup of finite index, itself have a polycyclic subgroup of finite index? The supplied registry labels this question open and attributes it to N. S. Chernikov (2006); that status is not independently verified here.
\subsection{Sources}
\begin{itemize}
\item Supplied local UnsolvedMath registry record, permanent ID \(701296\), exact source code \texttt{catalogue entry}, titled \emph{Kourovka Notebook Problem 16.105}. The record labels the problem open and was updated on 2026-10-16; this is supplied registry metadata, not independent literature verification.
\item \emph{The Kourovka Notebook (Unsolved Problems in Group Theory)}, Problem 16.105, proposed by N. S. Chernikov; the supplied record identifies the 16th issue (2006) and the 21st issue edition (2026), \texttt{arXiv:1401.0300v48}. The collection description names E. I. Khukhro and V. D. Mazurov as editors.
\item N. S. Chernikov, \emph{Ukrain. Math. J.} \textbf{32} (1980), 476--479, cited by the supplied statement for the comparison involving Chernikov factors. No article title or external link is supplied, and the cited result has not been independently checked here.
\end{itemize}
\end{document}
